\section{Introduction}
\label{sec:introduction}

Variational formulations turn a broad class of elliptic boundary-value problems into convex energy
minimization. Classical finite elements approximate the admissible space by locally supported
polynomials, and the relation between polynomial degree, mesh scale, and degrees of freedom is well
understood \citep{BrennerScott2008}. Neural and ridge-function methods instead construct global
nonlinear trial spaces, so that the approximation problem concerns the selection and
reoptimization of atoms rather than the design of a mesh. The deep Ritz method illustrates the
reach of variational neural solvers \citep{EYu2018}. The present work analyzes a deterministic
greedy selection rule for such nonlinear elliptic energies.

Orthogonal greedy approximation (OGA) has a sharp Hilbert-space theory. Under suitable source and
entropy conditions it produces optimal approximation rates
\citep{SiegelXuOGA2022,LiSiegelEntropy2024}. Chebyshev-type greedy methods extend the basic descent
argument to convex energies in Banach spaces \citep{TemlyakovCGA2013,TemlyakovGreedy2014}, and
biorthogonal or inexact variants show how selection and correction errors enter the energy recursion
\citep{DereventsovTemlyakov2022}. These results supply the approximation theory on which the analysis
below relies. Applying these results to a nonlinear partial differential equation (PDE)
approximated using atoms drawn from a finite candidate pool also requires
accounting for numerical quadrature, finite-dictionary approximation,
and inexact coefficient minimization. The three effects are distinct: quadrature
perturbs the evaluated score, a finite pool replaces the continuous supremum by a maximum over the
pool, and an active-space solver computes a minimizer to a prescribed accuracy.

The dictionary used below consists of shallow \(\relu^k\) ridge functions. Here
\(\relu^k(s)=[s]_+^k\), with \([s]_+=\max\{s,0\}\), and the domain is a bounded subset of
\(\R^d\). Greedy approximation with neural networks has been used successfully for PDEs
\citep{SiegelHongJinHaoXu2023}. A pool that is drawn once and then held fixed during the greedy
solve provides a way to formulate weak selection, by which we mean selection of an atom whose score is a
prescribed fraction of the best score available in the full dictionary:
\cite[Theorem~3.1]{XuXu2025} prove a convex-energy bound under Hilbert smoothness and
source/entropy assumptions, while their Theorems~5.1 and 5.2 quantify dictionary sizes that
realize weak selection. Their analysis assumes exact coefficient minimization over the active
span. The finite-pool inequality below separates the correction error from the score and coverage
errors. The universality of such constructions for nonlinear convex variational problems is
examined by \cite{BernaFalco2025}. Variation spaces and sharp entropy estimates for shallow
networks provide the source and width information used in the convergence estimates
\citep{SiegelXuVariation2023,SiegelXuSharp2024,MaoSiegelXu2026}.

The analysis has three parts. First, the one-step descent estimate separates pool coverage, score evaluation, radial stationarity, and coefficient-correction errors. Second, local Hessian transfer connects nonlinear greedy selection with weak OGA in a frozen Hilbert metric under relative-score, invariant-neighborhood, and entropy conditions. The entropy estimate uses the projected, unrenormalized restarted dictionary. Third, we estimate the projection error in the spaces generated by the
algorithm. These estimates account for the source remainder and the
observed orthogonal increments. For quartic energies, an exact identity
converts the projection estimates into nonlinear energy-error bounds. The numerical study evaluates the finite-trajectory estimates on the computed spaces.

For \(p\)-growth energies, the natural distance describes the geometry of the energy gap. The
nonlinear map
\[
  V(\xi)=|\xi|^{(p-2)/2}\xi
\]
defines the natural distance, as in the finite-element analysis of the \(p\)-Laplacian
\citep{BarrettLiu1993,DieningKreuzer2008,LeePark2025}. For the standard \(p\)-energies considered below, the energy gap is
comparable to the corresponding squared natural distance; hence an energy bound of order \(m^{-\alpha}\) gives a \(V\)-distance bound of order
\(m^{-\alpha/2}\) and, for \(p\geq2\), a \(W^{1,p}\)-error bound of order \(m^{-\alpha/p}\). This
conversion has the same form for every \(p\), although the energy exponent may depend on \(p\). We
prove the equivalence for the standard pure, regularized, and reaction \(p\)-energies and apply the
corresponding conversion to each of them.

Under the atomic source condition of \Cref{ass:atomic}, the directional smoothness estimate of \Cref{ass:directional}, and a uniform positive selection strength, the ideal continuous CGA has an \(O(m^{-1})\) energy bound. For finite-pool CGA, the one-step inequality quantifies separately the
effects of continuous-dictionary coverage, score-evaluation error,
radial active-space stationarity error, and coefficient-optimization error.

Local Hessian transfer combines Hessian coercivity, dictionary norm
comparison, and a relative score condition to obtain a weak-OGA
estimate in a frozen Hilbert metric.

The computed nested spaces also admit a finite-trajectory projection analysis. For a fixed finite symmetric dictionary and starting index \(m_0\), an approximate source representation with coefficient budget \(B_{\rm entry}\) and remainder bound \(\sigma\) gives
\[
 r_{m_0+j}^2\le \sigma^2+\frac{4B_{\rm entry}^2}{W_j},
 \qquad W_j=\sum_{m=m_0}^{m_0+j-1}\frac{\underline\theta_m^2}{\ell_m^2},
 \qquad W_j>0.
\]
Here \(r_m\) is the projection error, \(\ell_m\) is the innovation norm, and \(0\le\underline\theta_m\le\theta_m\) measures the selected score relative to the maximum frozen-dictionary score \(S_m^{\rm fin}\). For the one-dimensional \(p=4\) energies, an exact quartic expansion converts projection estimates into energy bounds. The analysis also gives an initial-value form of the source bound. Under the power-type condition \(\gamma_m^2\ge c_0r_m^{2+2/\alpha}\), with \(\gamma_m=\theta_m S_m^{\rm fin}/\ell_m\), the projection recurrence yields
\[
 r_{m_0+j}^2\le \left(r_{m_0}^{-2/\alpha}+c_0j/\alpha\right)^{-\alpha}.
\]
The quartic comparison then supplies the corresponding energy and metric bounds.

The numerical study covers eight cases: linear reaction--diffusion, cubic and hyperbolic-sine semilinear energies, and pure, regularized, and reaction \(p\)-energies. Five cases compare finite-pool CGA, random-feature models (RFM), and finite elements (FEM) at prescribed coefficient counts using a common evaluator. Random features follow the fixed-feature construction of \cite{RahimiRecht2007}; conditioning of linearized neural trial spaces is discussed by \cite{ChenChiEYang2022,MaoXuXu2026}. At the five reported endpoints, CGA has smaller Sobolev errors than the metric-valid RFM medians and smaller errors than the P3 FEM references in four cases.

The finite-trajectory study uses the one-dimensional pure and regularized \(p=4\) models. Source quantities and comparison constants are calibrated on \(N=8,\ldots,32\) and held fixed during continuation to \(N=64\). The approximate source is constructed at the starting index from the candidate and auxiliary reference dictionaries. The computations evaluate the source bound, the quartic energy identity, and the power-type projection bound with the prescribed benchmark \(\alpha_{\rm bench}=6\).

\Cref{sec:setting} introduces the variational setting, dictionaries, and algorithms. \Cref{sec:pgeometry} develops the energy families and natural-distance estimates. \Cref{sec:c5} presents finite-pool descent, local Hessian transfer, and finite-trajectory projection bounds. \Cref{sec:numerics} reports the CGA experiments, and \Cref{sec:results} compares the methods at prescribed coefficient counts. \Cref{sec:discussion} discusses the main findings.
